\chapter*{Abstract}
A charging location with only a few days of recorded demand offers little evidence about its daily pattern. This study asks how much a forecaster benefits from training at established locations and whether recent forecast errors indicate when to relax the influence of the source model. Using UrbanEV and CHARGED, we simulate cold start by limiting the initial target history to 7, 14 or 28 days. A gated recurrent unit uses the preceding 24 hours and known calendar features to predict hourly charging energy for the next day. We compare source initialization with target-only training, calibrating the initial training budget separately for each. We then compare fixed retention with online selection of the strength of a penalty on changes in the encoder representation. A separate experiment combines forecasts from two transferred models and tests the value of adapting their weights after observations arrive. On twelve additional UrbanEV targets, source initialization reduced mean absolute error by 29.80--33.35\% across the three history budgets. Across 21 historical combinations of dataset, city and history duration, confidence-guided retention selection yielded lower error than fixed retention in nine and lower error than selection by recent mean absolute error in eleven. Both the calibrated static mixture and the adaptive mixture improved on a single fixed-retention model in eighteen combinations. Adaptive fusion improved on the calibrated static blend in twelve. In a separate fourteen-day comparison on twelve UrbanEV targets, an equal-weight mixture came within 0.000007 kWh of the adaptive MAE. On 46 additional Development identifiers selected before evaluation at the 14-day budget, the estimated reductions in error from adaptive fusion were 0.67\% against single fixed retention and 0.36\% against static fusion when each dataset--city group received equal weight. Both uncertainty intervals included zero. The results support source initialization for the tested short histories. The benefit of adapting retention or forecast weights varied with the target group and the comparison method.

\vspace{0.5cm}
\noindent\textbf{Keywords:} electric vehicle charging, demand forecasting, transfer learning, cold start, online adaptation
